\documentclass[11pt,a4paper]{article}
\usepackage[utf8]{inputenc}
\usepackage[T1]{fontenc}
\usepackage{amsmath, amssymb, amsthm}
\usepackage{geometry}
\usepackage{hyperref}

\geometry{margin=1in}

\title{Game-Theoretic Formalization of the Decentralized Compliance Tribunal}
\author{Real World Assets (RWA) Tokenization Project}
\date{\today}

\newtheorem{theorem}{Theorem}
\newtheorem{lemma}{Lemma}
\newtheorem{definition}{Definition}

\begin{document}

\maketitle

\begin{abstract}
This document provides the formal game-theoretic analysis of the Compliance Tribunal utilized in our permissioned blockchain architecture. We demonstrate that under a specific configuration of reputational staking and slashing, truthful voting on compliance anomalies constitutes a strict Nash Equilibrium, establishing a robust Schelling Point without reliance on a centralized auditor.
\end{abstract}

\section{System Model}
Let $N$ be the set of randomly selected compliance auditors (peers) from regulated institutions (e.g., BNP Paribas, AMF), with $|N| = n$.
An AI agent flags an asset transfer as anomalous, triggering a Tribunal session. 
There exists an objective ground truth $\omega \in \{FRAUD, LEGITIMATE\}$.

Each auditor $i \in N$ observes evidence $E$ and must submit a vote $v_i \in \{FRAUD, LEGITIMATE\}$ using a cryptographic Commit-Reveal scheme to prevent cascading information cascades (free-riding).

\section{The Payoff Function}
Unlike public permissionless blockchains (e.g., Kleros, Augur) where auditors stake financial tokens, our permissioned environment uses a \textbf{Reputational Stake} $S \in \mathbb{R}^+$. A high reputation score is required to maintain node consensus privileges on the network.

The final decision $D$ is determined by a supermajority rule ($\tau = \frac{2}{3}n$).
If the votes are tallied and a supermajority is reached, $D$ is established. Otherwise, the asset remains frozen pending manual regulatory review.

The utility function $u_i(v_i, v_{-i})$ for an auditor $i$ is defined by the slashing fraction $f \in (0, 1]$ and a small reward $R$:
\begin{equation}
u_i(v_i, v_{-i}) = 
\begin{cases} 
S + R & \text{if } v_i = D \\
S - f \cdot S & \text{if } v_i \neq D \text{ and } D \text{ is reached}\\
S & \text{if no majority } D \text{ is reached}
\end{cases}
\end{equation}

\section{Incentive Compatibility and Schelling Point}
We assume that analyzing the evidence $E$ incurs a cognitive cost $C \ll S$. We also assume that auditors are rational and risk-averse regarding their institutional reputation $S$.

\begin{theorem}[Truthful Nash Equilibrium]
If all auditors share a common prior that the evidence $E$ clearly points to the ground truth $\omega$, and if the penalty $f \cdot S > R + C$, then the strategy profile where all auditors vote truthfully ($v_i = \omega \ \forall i$) is a strict Nash Equilibrium.
\end{theorem}
\begin{proof}
Assume the strategy profile $\sigma^*$ where $n-1$ auditors play $v_{-i} = \omega$.
If auditor $i$ also plays $v_i = \omega$, the supermajority $D = \omega$ is reached.
The utility is:
\[ u_i(\omega, \omega_{-i}) = S + R - C \]

If auditor $i$ deviates and plays $v_i \neq \omega$ (a malicious or lazy vote), the supermajority $D = \omega$ is still reached (assuming $n \geq 4$ and $\tau = 2/3$). The deviating auditor is slashed:
\[ u_i(\neq \omega, \omega_{-i}) = S - f \cdot S - C' \]
(where $C'$ is $0$ if they were lazy and didn't analyze $E$, or $C$ if they did).

For the truthful vote to be a strict best response, we require:
\[ S + R - C > S - f \cdot S \]
\[ f \cdot S > C - R \]

Since we set $f \cdot S > C + R$, the condition is strictly satisfied. Thus, no auditor has a rational incentive to deviate from the truthful consensus. This establishes the truthful outcome as a focal point (Schelling Point).
\end{proof}

\section{Resistance to Bribing Coalitions}
In a permissioned banking context, a malicious actor might attempt to bribe a coalition of size $k \geq \frac{n}{3}$ to prevent a supermajority against fraud, or $k \geq \frac{2n}{3}$ to force a false outcome.

\begin{lemma}[Coalition Unprofitability]
The cost of bribing a successful coalition $B_{cost}$ grows linearly with the reputational slash $f \cdot S$. 
\end{lemma}
Because $S$ represents institutional operational licenses (not just liquid capital), $S \to \infty$ from the perspective of an individual rogue trader. Therefore, the bribe required to convince $k$ institutions to risk their regulatory banking license is economically infeasible.

\section{Conclusion}
By adapting Schelling-point game theory to non-financial reputational stakes, the Compliance Tribunal provides mathematical guarantees for decentralized auditing without requiring a native cryptocurrency token.

\end{document}
